\section{Source Code and Adaptation}

The design was based on the folder \texttt{6DOF - CURSO/03-autopiloto}. The main script used as reference was \texttt{ProjRLocus.m}, which designs the attitude and acceleration autopilots by root locus using the short-period transfer functions generated by \texttt{FTransDin.m}.

For this work, the original script was copied to \texttt{ProjRLocus\_antiship.m}. The structure of the professor's code was preserved. The changes were limited to:

\begin{itemize}
    \item replacing the generic missile data by the antiship missile data from the previous 6DOF work;
    \item using the DATCOM aerodynamic matrix \texttt{M\_aed.mat} generated for the antiship missile;
    \item setting the design condition to Mach \(0.9\) and altitude \(10~\text{m}\);
    \item using representative mass/CG points along the burn;
    \item exporting the figures and the numerical requirement check.
\end{itemize}

The Simulink file \texttt{Autopiloto.slx} was kept in the code folder as the implementation model from the course material. The gain design itself follows the MATLAB root-locus script, which is the part of the professor's code intended for generating the autopilot tables.

\begin{lstlisting}[language=Matlab]
% Antiship missile design condition used in the assignment.
altitude = 10;
Mach = 0.9;
VetTmp = [0 D.TqB/2 D.TqB D.TqB + D.TqS/2];
\end{lstlisting}

The selected points cover launch mass, intermediate booster burn, booster burnout and a representative cruise point during the sustainer burn.
